\documentclass[a4paper,12pt]{article}

\usepackage[T2A]{fontenc}
\usepackage[utf8]{inputenc}
\usepackage[russian]{babel}

\usepackage{amsmath,amssymb,amsthm,mathtools}
\usepackage{helvet}
\renewcommand{\familydefault}{\sfdefault}

\usepackage{icomma}
\usepackage[a4paper,margin=2.15cm,headheight=16pt]{geometry}
\usepackage{microtype}

\usepackage{tikz}
\usepackage{tkz-euclide}
\usepackage{pgfplots}
\pgfplotsset{compat=1.18}

\usepackage{tabularx,booktabs,longtable}
\usepackage{enumitem}
\usepackage{hyperref}
\usepackage{parskip}
\usepackage{fancyhdr}
\usepackage{setspace}
\usepackage{xcolor}

% ------------------------------------------------------------
% Палитра
% ------------------------------------------------------------

\definecolor{accent}{HTML}{008080}
\definecolor{darkaccent}{HTML}{004D4D}
\definecolor{textgray}{HTML}{333333}
\definecolor{softgray}{HTML}{EEF3F3}
\definecolor{linegray}{HTML}{B7C7C7}

% ------------------------------------------------------------
% Общая типографика
% ------------------------------------------------------------

\onehalfspacing
\setlength{\parindent}{0pt}
\setlength{\parskip}{0.55em}

\setlist[itemize]{
    leftmargin=1.4em,
    itemsep=0.25em,
    topsep=0.25em
}

\setlist[enumerate]{
    leftmargin=1.7em,
    itemsep=0.45em,
    topsep=0.35em
}

\hypersetup{
    colorlinks=true,
    linkcolor=darkaccent,
    urlcolor=darkaccent
}

\pagestyle{fancy}
\fancyhf{}
\fancyhead[L]{\small\color{textgray}Векторы на координатной плоскости}
\fancyhead[R]{\small\color{textgray}Ксения Клыкова}
\fancyfoot[C]{\small\color{textgray}\thepage}
\renewcommand{\headrulewidth}{0.3pt}

\newcommand{\vect}[1]{\overrightarrow{#1}}
\newcommand{\dotprod}{\mathbin{\cdot}}

% ------------------------------------------------------------
% Документ
% ------------------------------------------------------------

\begin{document}

% ============================================================
% ТИТУЛЬНЫЙ ЛИСТ
% ============================================================

\begin{titlepage}
\thispagestyle{empty}

\begin{center}

\vspace*{2.3cm}

{\color{darkaccent}\Large ЧЕК-ЛИСТ}

\vspace{0.6cm}

{\Huge\bfseries
Векторы на координатной плоскости
}

\vspace{0.55cm}

{\Large
Координаты, длина, скалярное произведение,\\
угол между векторами и действия с векторами
}

\vspace{2.2cm}

{\large
Лёвин Артём Александрович\\[0.25cm]
эксклюзивно для Ксении Клыковой
}

\vspace{1.3cm}

{\large ЕГЭ по математике}

\vfill

{\color{textgray}\large \today}

\vspace{1.0cm}

\begin{tikzpicture}[x=1cm,y=1cm]
    \draw[
        color=accent,
        line width=1.2pt
    ]
    (0,0)
    .. controls (3,0.8) and (6,-0.8) .. (9,0)
    .. controls (11,0.45) and (13,0.45) .. (15,0);
\end{tikzpicture}

\vspace{0.7cm}

\end{center}

\end{titlepage}

% ============================================================
% ОСНОВНОЙ ЧЕК-ЛИСТ
% ============================================================

\section*{\color{darkaccent}1. Вектор и его координаты}

\textbf{Вектор} --- направленный отрезок. У него есть начало и конец.

Если точка \(A\) является началом, а точка \(B\) --- концом, то вектор обозначают
\[
\vect{AB}.
\]

Пусть
\[
A(x_A;y_A),
\qquad
B(x_B;y_B).
\]

Тогда координаты вектора \(\vect{AB}\) вычисляются по формуле
\[
\boxed{
\vect{AB}=(x_B-x_A;\;y_B-y_A)
}
\]

Главное правило:

\[
\boxed{\text{координата конца} - \text{координата начала}}
\]

Это правило применяется отдельно к каждой координате.

\subsection*{Пример}

Пусть
\[
A(2;3),
\qquad
B(8;6).
\]

Тогда
\[
\vect{AB}=(8-2;\;6-3)=(6;3).
\]

Следовательно, вектор смещает точку:

\begin{itemize}
    \item на \(6\) единиц вправо;
    \item на \(3\) единицы вверх.
\end{itemize}

\subsection*{Координаты по смещению}

Если вектор изображён на координатной сетке, его координаты часто можно определить без вычислений.

\[
\begin{array}{c|c}
\text{Смещение} & \text{Знак координаты} \\
\midrule
\text{вправо} & x>0 \\
\text{влево} & x<0 \\
\text{вверх} & y>0 \\
\text{вниз} & y<0
\end{array}
\]

Например:

\[
\text{на 4 вправо и на 2 вниз}
\quad\Longrightarrow\quad
(4;-2).
\]

\textbf{Важно.}
Координаты вектора показывают именно \emph{смещение}, а не положение его начала или конца.

% ------------------------------------------------------------

\section*{\color{darkaccent}2. Длина вектора}

Пусть
\[
\vec a=(x;y).
\]

Длина вектора обозначается
\[
|\vec a|
\]
и вычисляется по формуле

\[
\boxed{
|\vec a|=\sqrt{x^2+y^2}
}
\]

Формула является непосредственным следствием теоремы Пифагора.

\subsection*{Пример}

Для вектора
\[
\vec a=(6;3)
\]
имеем

\[
|\vec a|
=
\sqrt{6^2+3^2}
=
\sqrt{36+9}
=
\sqrt{45}
=
3\sqrt5.
\]

\textbf{Типичная ошибка:}
подставлять в формулу координаты точек \(A\) и \(B\).

Сначала, если необходимо, найдите
\[
\vect{AB}=(x_B-x_A;y_B-y_A),
\]
и только затем вычисляйте
\[
|\vect{AB}|.
\]

% ------------------------------------------------------------

\section*{\color{darkaccent}3. Скалярное произведение через координаты}

Пусть
\[
\vec a=(x_a;y_a),
\qquad
\vec b=(x_b;y_b).
\]

Скалярное произведение вычисляется по формуле

\[
\boxed{
\vec a\dotprod\vec b
=
x_a x_b+y_a y_b
}
\]

То есть:

\begin{enumerate}
    \item перемножьте соответствующие \(x\)-координаты;
    \item перемножьте соответствующие \(y\)-координаты;
    \item сложите полученные произведения.
\end{enumerate}

\subsection*{Пример}

Пусть
\[
\vec a=(1;2),
\qquad
\vec b=(3;4).
\]

Тогда

\[
\vec a\dotprod\vec b
=
1\cdot3+2\cdot4
=
3+8
=
11.
\]

\textbf{Следите за знаками.}

Например, если
\[
\vec k=(-1;3),
\qquad
\vec p=(2;-5),
\]
то

\[
\vec k\dotprod\vec p
=
(-1)\cdot2+3\cdot(-5)
=
-2-15
=
-17.
\]

% ------------------------------------------------------------

\section*{\color{darkaccent}4. Скалярное произведение через угол}

Если \(\varphi\) --- угол между ненулевыми векторами
\(\vec a\) и \(\vec b\), то

\[
\boxed{
\vec a\dotprod\vec b
=
|\vec a|\,|\vec b|\cos\varphi
}
\]

Сопоставляя две формулы скалярного произведения, получаем

\[
x_a x_b+y_a y_b
=
|\vec a|\,|\vec b|\cos\varphi.
\]

Отсюда

\[
\boxed{
\cos\varphi
=
\frac{x_a x_b+y_a y_b}
{\sqrt{x_a^2+y_a^2}\,
 \sqrt{x_b^2+y_b^2}}
}
\]

Формула применима при

\[
|\vec a|\ne0,
\qquad
|\vec b|\ne0.
\]

То есть оба вектора должны быть \textbf{ненулевыми}.

% ------------------------------------------------------------

\section*{\color{darkaccent}5. Алгоритм нахождения угла между векторами}

Если требуется найти угол между векторами, действуйте последовательно.

\begin{enumerate}
    \item Найдите координаты каждого вектора, если они не заданы явно.
    \item Вычислите скалярное произведение
    \[
    \vec a\dotprod\vec b.
    \]
    \item Вычислите длины
    \[
    |\vec a|,
    \qquad
    |\vec b|.
    \]
    \item Найдите
    \[
    \cos\varphi
    =
    \frac{\vec a\dotprod\vec b}
    {|\vec a|\,|\vec b|}.
    \]
    \item Если значение табличное, определите сам угол.
\end{enumerate}

\subsection*{Табличные значения косинуса}

\begin{center}
\renewcommand{\arraystretch}{1.3}
\begin{tabular}{ccccc}
\toprule
\(\varphi\) &
\(0^\circ\) &
\(30^\circ\) &
\(45^\circ\) &
\(60^\circ\) \\
\midrule
\(\cos\varphi\) &
\(1\) &
\(\dfrac{\sqrt3}{2}\) &
\(\dfrac{\sqrt2}{2}\) &
\(\dfrac12\) \\
\bottomrule
\end{tabular}
\end{center}

Кроме того,
\[
\cos90^\circ=0.
\]

\subsection*{Пример}

Пусть
\[
\vec a=(1;2),
\qquad
\vec b=(3;4).
\]

Тогда

\[
\vec a\dotprod\vec b=11,
\]

\[
|\vec a|=\sqrt5,
\qquad
|\vec b|=5.
\]

Следовательно,

\[
\cos\varphi
=
\frac{11}{5\sqrt5}.
\]

Если значение косинуса не является привычным табличным значением, в задачах без требования найти численную меру угла ответ может быть оставлен в таком виде.

% ------------------------------------------------------------

\section*{\color{darkaccent}6. Что можно узнать по знаку скалярного произведения}

Поскольку

\[
\vec a\dotprod\vec b
=
|\vec a|\,|\vec b|\cos\varphi,
\]

для ненулевых векторов знак скалярного произведения совпадает со знаком \(\cos\varphi\).

\begin{center}
\renewcommand{\arraystretch}{1.35}
\begin{tabularx}{\linewidth}{cX}
\toprule
Условие & Вывод \\
\midrule
\(\vec a\dotprod\vec b>0\)
&
угол между векторами острый
\\

\(\vec a\dotprod\vec b=0\)
&
векторы перпендикулярны
\\

\(\vec a\dotprod\vec b<0\)
&
угол между векторами тупой
\\
\bottomrule
\end{tabularx}
\end{center}

Например,

\[
(-1;3)\dotprod(2;-5)=-17<0,
\]

поэтому угол между этими векторами является тупым.

% ------------------------------------------------------------

\section*{\color{darkaccent}7. Сложение и вычитание векторов}

Пусть

\[
\vec a=(x_a;y_a),
\qquad
\vec b=(x_b;y_b).
\]

Тогда

\[
\boxed{
\vec a+\vec b
=
(x_a+x_b;\;y_a+y_b)
}
\]

и

\[
\boxed{
\vec a-\vec b
=
(x_a-x_b;\;y_a-y_b)
}
\]

Координаты складываются и вычитаются \textbf{покомпонентно}.

\subsection*{Пример}

Если

\[
\vec a=(1;2),
\qquad
\vec b=(3;4),
\]

то

\[
\vec a+\vec b=(4;6),
\]

\[
\vec a-\vec b=(-2;-2).
\]

% ------------------------------------------------------------

\section*{\color{darkaccent}8. Умножение вектора на число}

Если

\[
\vec a=(x;y)
\]

и \(\lambda\) --- число, то

\[
\boxed{
\lambda\vec a=(\lambda x;\lambda y)
}
\]

То есть число умножается на \textbf{каждую} координату.

Например,

\[
-2(3;-4)=(-6;8).
\]

При отрицательном коэффициенте особенно внимательно следите за знаками.

% ------------------------------------------------------------

\section*{\color{darkaccent}9. Линейные комбинации векторов}

Вектор может быть задан не своими координатами, а через другие векторы.

Например,

\[
\vec d=2\vec a-3\vec b,
\]

где

\[
\vec a=(1;2),
\qquad
\vec b=(3;4).
\]

Лучше выполнять вычисления по шагам.

\[
2\vec a=(2;4),
\]

\[
3\vec b=(9;12),
\]

поэтому

\[
\vec d
=
2\vec a-3\vec b
=
(2;4)-(9;12)
=
(-7;-8).
\]

Теперь можно найти длину:

\[
|\vec d|
=
\sqrt{(-7)^2+(-8)^2}
=
\sqrt{49+64}
=
\sqrt{113}.
\]

\textbf{Методический принцип:}
не пытайтесь выполнять длинную цепочку вычислений мысленно.
Сначала найдите координаты промежуточных векторов, затем выполняйте сложение или вычитание.

% ------------------------------------------------------------

\section*{\color{darkaccent}10. Ещё один типовой пример}

Пусть

\[
\vec t=(-2;3),
\qquad
\vec q=(0;5),
\]

а

\[
\vec n=-2\vec t+3\vec q.
\]

Вычислим по шагам:

\[
-2\vec t
=
-2(-2;3)
=
(4;-6),
\]

\[
3\vec q
=
3(0;5)
=
(0;15).
\]

Следовательно,

\[
\vec n
=
(4;-6)+(0;15)
=
(4;9).
\]

Поэтому

\[
|\vec n|
=
\sqrt{4^2+9^2}
=
\sqrt{97}.
\]

% ------------------------------------------------------------

\section*{\color{darkaccent}11. Работа с корнями в задачах на угол}

После подстановки в формулу косинуса часто требуется упростить выражение.

Например,

\[
\cos\varphi
=
\frac{5}{\sqrt5\sqrt{10}}.
\]

Так как

\[
\sqrt{10}=\sqrt5\sqrt2,
\]

получаем

\[
\cos\varphi
=
\frac{5}{\sqrt5\cdot\sqrt5\cdot\sqrt2}
=
\frac{5}{5\sqrt2}
=
\frac{1}{\sqrt2}.
\]

Рационализируем знаменатель:

\[
\frac{1}{\sqrt2}
=
\frac{\sqrt2}{2}.
\]

Следовательно,

\[
\cos\varphi=\frac{\sqrt2}{2}
\quad\Longrightarrow\quad
\boxed{\varphi=45^\circ}.
\]

% ------------------------------------------------------------

\section*{\color{darkaccent}12. Основной алгоритм решения задач}

Если дана задача на векторы, удобно придерживаться одного порядка действий.

\begin{enumerate}
    \item Определите, какие векторы участвуют в задаче.
    \item Если вектор задан двумя точками, найдите его координаты:
    \[
    \text{конец}-\text{начало}.
    \]
    \item Если вектор задан выражением типа
    \[
    2\vec a-3\vec b,
    \]
    сначала найдите его координаты.
    \item Только после этого применяйте нужную формулу:
    \begin{itemize}
        \item длина;
        \item скалярное произведение;
        \item косинус угла;
        \item сумма или разность.
    \end{itemize}
    \item Проверьте знаки.
    \item Аккуратно упростите корни и дроби.
    \item Оцените ответ по смыслу.
\end{enumerate}

% ------------------------------------------------------------

\section*{\color{darkaccent}13. Быстрая проверка ответа}

Перед завершением задачи полезно проверить несколько вещей.

\begin{itemize}
    \item Длина вектора не может быть отрицательной.
    \item Для ненулевых векторов
    \[
    -1\le\cos\varphi\le1.
    \]
    \item Если получилось
    \[
    |\cos\varphi|>1,
    \]
    в вычислениях есть ошибка.
    \item Если скалярное произведение отрицательно, угол должен быть тупым.
    \item Если скалярное произведение равно нулю, угол равен
    \[
    90^\circ.
    \]
    \item В координатах \(\vect{AB}\) всегда вычитается
    \[
    B-A,
    \]
    а не наоборот.
\end{itemize}

% ------------------------------------------------------------

\section*{\color{darkaccent}14. Типичные ошибки}

\begin{enumerate}
    \item \textbf{Перепутать начало и конец вектора.}

    Для \(\vect{AB}\):
    \[
    (x_B-x_A;y_B-y_A).
    \]

    \item \textbf{Потерять знак минус.}

    Например,
    \[
    3\cdot(-5)=-15.
    \]

    \item \textbf{При умножении вектора на число изменить только одну координату.}

    Неверно:
    \[
    3(2;-1)=(6;-1).
    \]

    Верно:
    \[
    3(2;-1)=(6;-3).
    \]

    \item \textbf{Перепутать сложение и вычитание при работе с линейной комбинацией.}

    Если
    \[
    \vec n=-2\vec t+3\vec q,
    \]
    после нахождения \(-2\vec t\) и \(3\vec q\) эти два вектора нужно
    \textbf{сложить}.

    \item \textbf{Забыть скобки около отрицательных координат.}

    Лучше записывать
    \[
    (-2)^2,
    \]
    а не
    \[
    -2^2.
    \]

    Действительно,
    \[
    (-2)^2=4,
    \qquad
    -2^2=-4.
    \]

    \item \textbf{Смешать координаты точек и координаты вектора.}

    Формула
    \[
    |\vec a|=\sqrt{x^2+y^2}
    \]
    использует координаты \textbf{самого вектора}.

    \item \textbf{Забыть, что формула угла требует ненулевых векторов.}

    В знаменателе находятся их длины:
    \[
    |\vec a|\,|\vec b|\ne0.
    \]
\end{enumerate}

% ------------------------------------------------------------

\section*{\color{darkaccent}15. Формулы, которые нужно знать}

\[
\boxed{
\vect{AB}
=
(x_B-x_A;\;y_B-y_A)
}
\]

\[
\boxed{
|\vec a|
=
\sqrt{x_a^2+y_a^2}
}
\]

\[
\boxed{
\vec a\dotprod\vec b
=
x_a x_b+y_a y_b
}
\]

\[
\boxed{
\vec a\dotprod\vec b
=
|\vec a|\,|\vec b|\cos\varphi
}
\]

\[
\boxed{
\cos\varphi
=
\frac{x_a x_b+y_a y_b}
{\sqrt{x_a^2+y_a^2}\sqrt{x_b^2+y_b^2}}
}
\]

\[
\boxed{
\vec a+\vec b
=
(x_a+x_b;\;y_a+y_b)
}
\]

\[
\boxed{
\vec a-\vec b
=
(x_a-x_b;\;y_a-y_b)
}
\]

\[
\boxed{
\lambda\vec a
=
(\lambda x_a;\lambda y_a)
}
\]

\clearpage

% ============================================================
% ТРЕНИРОВОЧНЫЙ БЛОК
% ============================================================

\section*{\color{darkaccent}Тренировочный блок}

Выполняйте задания последовательно. Особое внимание уделяйте знакам и порядку координат.

\subsection*{Средний уровень}

\begin{enumerate}[label=\textbf{\arabic*.}]

\item
Даны точки
\[
A(-2;3),
\qquad
B(5;-1).
\]
Найдите координаты вектора \(\vect{AB}\).

\item
Дан вектор
\[
\vec a=(-6;8).
\]
Найдите его длину.

\item
Даны векторы
\[
\vec a=(2;-3),
\qquad
\vec b=(-4;5).
\]
Найдите их скалярное произведение.

\end{enumerate}

\subsection*{Выше среднего}

\begin{enumerate}[label=\textbf{\arabic*.},start=4]

\item
Даны векторы
\[
\vec a=(3;-2),
\qquad
\vec b=(-1;4).
\]
Найдите координаты вектора
\[
\vec c=2\vec a-\vec b.
\]

\item
Даны
\[
\vec p=(-2;1),
\qquad
\vec q=(3;-4).
\]
Найдите длину вектора
\[
\vec r=-3\vec p+2\vec q.
\]

\item
Даны векторы
\[
\vec a=(1;0),
\qquad
\vec b=(1;\sqrt3).
\]
Найдите угол между ними.

\item
Даны векторы
\[
\vec a=(2;1),
\qquad
\vec b=(1;-2).
\]
Докажите вычислением скалярного произведения, что они перпендикулярны.

\item
Даны точки
\[
A(-1;2),
\qquad
B(2;3),
\qquad
C(4;-1),
\qquad
D(2;-2).
\]
Найдите
\[
\vect{AB}\dotprod\vect{CD}.
\]

\item
Даны точки
\[
A(1;1),
\qquad
B(3;2),
\qquad
C(-2;4),
\qquad
D(1;3).
\]
Найдите косинус угла между векторами
\[
\vect{AB}
\quad\text{и}\quad
\vect{CD}.
\]

\end{enumerate}

\subsection*{Сложный уровень}

\begin{enumerate}[label=\textbf{\arabic*.},start=10]

\item
Даны векторы
\[
\vec a=(1;2),
\qquad
\vec b=(2;-1).
\]

Вектор \(\vec c\) определяется равенством
\[
\vec c=3\vec a-2\vec b,
\]
а вектор \(\vec d\) --- равенством
\[
\vec d=\vec a+\vec b.
\]

Найдите угол между векторами \(\vec c\) и \(\vec d\).

\end{enumerate}

\clearpage

% ============================================================
% ОТВЕТЫ
% ============================================================

\section*{\color{darkaccent}Ответы к тренировочному блоку}

\begin{table}[ht]
\centering
\caption{Ответы}
\renewcommand{\arraystretch}{1.4}

\begin{tabularx}{\linewidth}{cX}
\toprule
\textbf{№} & \textbf{Ответ} \\
\midrule

1 &
\[
\vect{AB}=(7;-4)
\]
\\

2 &
\[
|\vec a|=10
\]
\\

3 &
\[
\vec a\dotprod\vec b=-23
\]
\\

4 &
\[
\vec c=(7;-8)
\]
\\

5 &
\[
\vec r=(12;-11),
\qquad
|\vec r|=\sqrt{265}
\]
\\

6 &
\[
\cos\varphi=\frac12,
\qquad
\varphi=60^\circ
\]
\\

7 &
\[
\vec a\dotprod\vec b
=
2\cdot1+1\cdot(-2)
=
0,
\qquad
\varphi=90^\circ
\]
\\

8 &
\[
\vect{AB}=(3;1),
\qquad
\vect{CD}=(-2;-1),
\qquad
\vect{AB}\dotprod\vect{CD}=-7
\]
\\

9 &
\[
\vect{AB}=(2;1),
\qquad
\vect{CD}=(3;-1),
\qquad
\cos\varphi=\frac{1}{\sqrt2}
=
\frac{\sqrt2}{2}
\]
\\

10 &
\[
\vec c=(-1;8),
\qquad
\vec d=(3;1),
\]
\[
\cos\varphi
=
\frac{5}{\sqrt{65}\sqrt{10}}
=
\frac{1}{\sqrt{26}}
=
\frac{\sqrt{26}}{26},
\]
\[
\varphi
=
\arccos\!\left(\frac{1}{\sqrt{26}}\right)
\]
\\

\bottomrule
\end{tabularx}

\end{table}

\vfill

\begin{center}
{\color{darkaccent}\rule{0.55\linewidth}{0.6pt}}

\vspace{0.5cm}

\textbf{Главная идея}

Сначала определите координаты нужного вектора,\\
затем спокойно применяйте соответствующую формулу.

\end{center}

\end{document}